\chapter{ART, Ageing and the Family}

\section{Assisted Reproductive Technology (ART)}

\begin{KeyConcept}
\begin{itemize}
\item \textbf{ART} (the 'genetic revolution' in the family) includes \textbf{IVF, intracytoplasmic sperm injection, cryopreservation of eggs/sperm, and artificial insemination}. It was created to cure infertility but expanded its scope.
\item \textbf{Criticisms} (legal + social): (1) \textbf{Inequality} --- ART is costly, so the rich can produce children with 'socially acceptable' qualities and genetic modification, while the poor cannot (like educational inequality once created caste); (2) \textbf{surrogates are from poor backgrounds}, have \textbf{no legal rights} (only a contract), and exchange their bodily privacy for unregulated money --- India and Nepal are the \textbf{surrogacy hubs} (controlled by mafias; the government moved to ban commercial surrogacy).
\item ART is also \textbf{fundamentally conservative}: it produces \textbf{biologically related children} (unlike adoption), so it reproduces the same biological/racial idea of family; the \textbf{sperm donor has no rights}, and courts inconsistently grant divorced fathers rights.
\item But ART also \textbf{enables families} for gay, lesbian, single and \textbf{post-menopausal women}, thereby \textbf{assaulting the traditional image of heteronormativity}.
\item Three ways ART disturbs the traditional family: it creates \textbf{social parenthood} (not biological --- artificial insemination), it \textbf{destroys the family/market binary} (commercialisation of reproduction), and it \textbf{privatises parenthood} (contract in place of social bond).
\item Biological vs social parenthood is legally unresolved: in the \textbf{N. D. Tiwari} case, DNA testing proved paternity and gave the child \emph{some} property rights --- but \textbf{not equal} (biological) rights, unlike a socially-adopted child who has full rights. \textbf{Single parents using ART are a legal 'grey area'} in India. The \textbf{Nepal 2015 earthquake} stranded surrogate mothers (many under contract with Israeli parents, for whom Israel sent special assistance) --- exposing the unregulated, exploitative market.
\item \textbf{Emancipation for upper-class women} (pregnancy without being pregnant) rests on the \textbf{exploitation of lower-class women}. \textbf{Ogburn's 'cultural lag'}: material development has outpaced cultural/legal development (India has no ART law --- it uses archaic IPC/CrPC provisions; the West developed norms through debate --- e.g. the 'three-parent baby') --- a lag that produces \textbf{anomie} (the old values destroyed but new ones not yet established).
\end{itemize}
\end{KeyConcept}

\begin{KeyConcept}
\begin{itemize}
\item \textbf{Solution --- land rights}: land is the most basic productive asset (Kerala's successful land reform vs Bihar's failure --- which is why Kerala's human development is far ahead); the \textbf{RTE} example shows how giving rights (education) empowered even rural women to protest. \textbf{David Harvey} (critic of neoliberalism): land is the last sphere the market could not fully control (family/honour/subsistence), which is why the farm laws were resisted. \textbf{Neoliberalism} = the retreat of the state from the market; it created a \textbf{Gramscian hegemony} (the 'common sense' that every young person must become an engineer).
\end{itemize}
\end{KeyConcept}

\begin{QuickRevision}
\begin{itemize}
\item ART = IVF, sperm injection, cryopreservation, artificial insemination.
\item Criticisms: inequality (rich $\rightarrow$ genetic modification), surrogates (poor, no rights, hubs India/Nepal, mafias), conservative (biological children), but enables gay/single/post-menopausal families.
\item Disturbs family: social parenthood, family/market binary broken, privatised parenthood; upper-class emancipation vs lower-class exploitation; Ogburn's cultural lag.
\item Solution: land rights (Kerala vs Bihar; RTE); David Harvey (land = last uncontrolled sphere); neoliberalism (retreat of state; Gramsci's hegemony).
\end{itemize}
\end{QuickRevision}

\section{Ageing and Old-Age Problems}

\begin{KeyConcept}
\begin{itemize}
\item \textbf{Ageing} (60+) --- most of the aged live in rural India. Problems:
\item \textbf{Modernisation} --- isolation and alienation (dual-working families; children migrate to cities), reflected in the nature of care.
\item \textbf{Social welfare under neoliberalism} --- pension funds invested in the volatile market (\textbf{casino capitalism}; the 2007--08 crisis wiped out pensions); hospitals, schools, parks privatised (gated communities).
\item \textbf{Widowhood} --- women live longer than men and, having not worked and holding no assets, face no social security and family support; rural widows are even stereotyped as \textbf{witches}.
\item \textbf{New diseases} (Alzheimer's, cancer) $\rightarrow$ \textbf{geriatric care} is emerging but weak (the aged are not 'part of the economy' or politics). \textbf{Old-age homes} give care and community but alienation returns when visitors leave.
\item \textbf{Social exclusion} --- the \textbf{digital divide} (biometric verification for pensions, online banking), financial exclusion, and exclusion from the labour market (the demographic dividend's cheap young labour makes the old's experience unwanted); their \textbf{roles} (socialising grandchildren) have been taken by schools and hospitals.
\item \textbf{Zygmunt Bauman} ('liquid modernity'): the \textbf{place of death shifted from the household to the hospital}; \textbf{Ivan Illich's 'medicalization'}: medicine has taken over the roles of law and religion (including death, cremation, palliative care and even \textbf{euthanasia}). People also \textbf{resist} --- opting for alternate medicine and not chemotherapy, to preserve their relationships/identity.
\item \textbf{Kinship has changed}: from blood/adoption to \textbf{fictive kinship} ('uncle', 'bhaiya') --- a move \textbf{from kinship to friendship}.
\end{itemize}
\end{KeyConcept}

\begin{QuickRevision}
\begin{itemize}
\item Ageing (60+; rural): isolation (dual-working families), neoliberal welfare (casino capitalism, privatisation), widowhood (witches), new diseases $\rightarrow$ geriatric care, old-age homes, digital divide/biometric exclusion, roles taken by schools/hospitals.
\item Bauman (death $\rightarrow$ hospital), Illich (medicalization $\rightarrow$ medicine took law/religion's roles); fictive kinship (kinship $\rightarrow$ friendship).
\end{itemize}
\end{QuickRevision}

\section{The Family: Murdock, Parsons and the Symmetrical Family}

\begin{KeyConcept}
\begin{itemize}
\item \textbf{George Peter Murdock} (functionalist): the \textbf{family} is a social group of \textbf{common residence, economic cooperation and reproduction}, including adults of both sexes in a \textbf{socially approved sexual relationship} and their (biological or adopted) children --- performing four functions (sexual, reproductive, economic, educational). E.g. the \textbf{Banaro tribe} of New Guinea (the husband does not have sex with the wife until a friend produces the children).
\item \textbf{Talcott Parsons}: the family's role is \textbf{primary socialisation and the internalisation of cultural norms}.
\item \textbf{Michael Young and Peter Willmott ('The Symmetrical Family')} --- three stages: (1) \textbf{pre-industrial} (the household is the production unit, power skewed to men); (2) \textbf{early industrial} (production moves to the public sphere, male mortality high $\rightarrow$ women formed an \textbf{'informal trade union'} of kin/neighbours for insurance, so the conjugal bond weakened); (3) \textbf{the symmetrical family} (isolated nuclear household, \textbf{conjugal roles shared}, strong bond). Reasons for symmetry: \textbf{reduction of kinship role} (real wages, the welfare state), \textbf{geographical mobility}, \textbf{fewer children} (6 $\rightarrow$ 2), and the \textbf{home becoming attractive} (higher income, parties at home).
\item \textbf{Family vs household}: policy has shifted from 'family' to 'household' (MGNREGA is per household) --- 'family' is value-laden and conservative, while 'household' is inclusive of live-in relationships and emerging family forms.
\end{itemize}
\end{KeyConcept}

\begin{KeyConcept}
\begin{itemize}
\item \textbf{Feminist critique}: \textbf{Ann Oakley} (London) found conjugal equality is a myth --- women still do most housework even when working; the real problem is the \textbf{sexual division of labour}. \textbf{Sylvia Walby}: patriarchy has not disappeared but has taken \textbf{new forms}, moving from the private to the public sphere.
\item \textbf{Marital breakdown}: \textbf{divorce}, \textbf{separation}, and \textbf{empty-shell marriages} (living together apart).
\end{itemize}
\end{KeyConcept}

\begin{QuickRevision}
\begin{itemize}
\item Murdock (4 functions; Banaro tribe), Parsons (primary socialisation), Young \& Willmott (3 stages; symmetrical family).
\item Family vs household (MGNREGA; live-in); feminist critique (Oakley --- conjugal inequality, sexual division of labour; Walby --- patriarchy new forms); marital breakdown (divorce, separation, empty-shell).
\end{itemize}
\end{QuickRevision}
